\begin{abstract}
Supporting exercise decisions requires more than recording distance and pace: users must also interpret environmental conditions, available resources, and incomplete information. This paper presents AI Fitness Coach, a proposed context-aware artificial intelligence of things (AIoT) architecture that combines smartphone sensing, geographic information, weather data, and user input without requiring a heart-rate wearable. A domain agent uses a large language model for intent recognition and task planning and coordinates domain tools through the Model Context Protocol. The proposed workflow considers accessibility and detour costs when recommending points of interest, and combines food-image interpretation, user confirmation, and nutrition database retrieval for dietary assessment. Text, voice, and a three-dimensional avatar are planned interaction modes. A separate data-management prototype uses wallet authentication, AES-256-GCM encryption, and SHA-256 digests; an IOTA EVM anchoring interface is planned, but real chain writes remain disabled. Source development records report 56 passing local and simulated tests across security, assistant integration, avatar, and audio modules. These tests do not establish end-to-end performance or clinical benefit. We present a comparative evaluation plan for recommendation quality, nutritional estimation, interaction, robustness, and integrity verification. Exercise safety and long-term adherence benefits remain to be validated.
\end{abstract}
\begin{IEEEkeywords}
AIoT, intelligent agents, context awareness, Model Context Protocol, exercise, data integrity
\end{IEEEkeywords}
